\begin{abstract}
We formalize the averaging dynamics on a finite graph, in which every node replaces its value by
the average of its neighbours' values: the degree-weighted sum is conserved, the values stay in
the initial range, and on a connected graph with an odd closed walk all values converge to the
degree-weighted average of the initial values; on a connected bipartite graph they need not
converge. On a well-clustered regular graph started from random signs, the sign of the last change
recovers the two clusters exactly (Becchetti, Clementi, Natale, Pasquale, Trevisan, SODA 2017).
\end{abstract}

\section{The dynamics}
\begin{definition}[Averaging]\label{def:avg}
\lean{Averaging.avgStep,Averaging.avgIter,Averaging.degAvg}\leanok
$x_{t+1}(v)=\frac1{\deg v}\sum_{u\sim v}x_t(u)$, and $\bar x=\sum_v\deg(v)\,x_0(v)/\sum_v\deg(v)$.
\end{definition}
\begin{lemma}[Conservation and range]\label{lem:basic}
\lean{Averaging.degree_weighted_sum_step,Averaging.avgStep_le,Averaging.le_avgStep,Averaging.degree_weighted_sum_iter}\leanok\uses{def:avg}
$\sum_v\deg(v)\,x_{t}(v)$ does not depend on $t$; without isolated nodes, a round never leaves an
interval containing all current values.
\end{lemma}
\begin{proof}\leanok
Double counting over ordered adjacent pairs; an average of values in an interval stays in it.
\end{proof}

\section{Convergence}
\begin{lemma}[Walk weights]\label{lem:walks}
\lean{Averaging.transW,Averaging.transW_nonneg,Averaging.transW_sum,Averaging.avgIter_eq_sum_transW,Averaging.transW_pos_of_walk,Averaging.exists_uniform_walk_length}\leanok\uses{def:avg}
$x_t(v)=\sum_u W_t(v,u)\,x_0(u)$ with $W_t$ stochastic and $W_t(v,u)>0$ when a walk of length $t$
joins $v$ and $u$. On a connected graph with an odd closed walk, some length $L$ joins every pair.
\end{lemma}
\begin{proof}\leanok
Induction on $t$; for $L$, fix the parity with the odd closed walk and pad with back-and-forth steps.
\end{proof}
\begin{theorem}[Convergence to the degree-weighted average]\label{thm:conv}
\lean{Averaging.fMax_antitone,Averaging.fMin_monotone,Averaging.range_avgIter_le,Averaging.gap_iter,Averaging.tendsto_degAvg}\leanok\uses{lem:basic,lem:walks}
On a connected graph with an odd closed walk, $x_t(v)\to\bar x$ for every $v$.
\end{theorem}
\begin{proof}\leanok
All entries of $W_L$ are at least some $\delta>0$, so $\max-\min$ contracts by $1-|V|\delta$ every
$L$ rounds (Doeblin); the common limit is $\bar x$ by conservation.
\end{proof}
\begin{theorem}[Bipartite graphs]\label{thm:bip}
\lean{Averaging.not_tendsto_of_colorable}\leanok\uses{def:avg}
On a connected bipartite graph with at least two nodes, the values $\pm1$ of a proper
2-colouring change sign at every round, so they converge nowhere.
\end{theorem}
\begin{proof}\leanok
All neighbours of a node have the other colour.
\end{proof}

\section{Strong reconstruction}
\begin{definition}[Clustered regular graph]\label{def:clustered}
\lean{Averaging.IsBalancedPartition,Averaging.IsClusteredRegular,Averaging.clusterIndicator,Averaging.transitionMatrix,Averaging.maxAbsOtherEigenvalue,Averaging.color,Averaging.IsStrongReconstruction,Averaging.reconstructionTime}\leanok\uses{def:avg}
$G$ is $d$-regular on two clusters $V_1,V_2$ of $n$ nodes, each node having $b$ neighbours in the
other cluster; $\chi=\mathbb 1_{V_1}-\mathbb 1_{V_2}$, $P=A/d$ with eigenvalues
$\lambda_1\ge\lambda_2\ge\dots$, $\lambda=\max_{i\ge3}|\lambda_i|$. Node $v$ is blue at round $t$
if $x_t(v)\ge x_{t-1}(v)$; a coloring is a strong reconstruction if no color appears in both
clusters. $T(n,\delta)=\lceil\log(4n^3)/\log(1+\delta)\rceil+1$.
\end{definition}
\begin{lemma}[Lemma C.1]\label{lem:c1}
\lean{Averaging.avgIter_eq_transitionMatrix_pow_mulVec,Averaging.transitionMatrix_mulVec_clusterIndicator,Averaging.abs_avgIter_sub_le}\leanok\uses{def:clustered}
If $\lambda<1-2b/d$ and $x\in\{\pm1\}^{2n}$, then
$|x_t(v)-\alpha_1-\alpha_2(1-2b/d)^t\chi(v)|\le\lambda^t\sqrt{2n}$ with
$\alpha_1=\langle x,\mathbb 1\rangle/2n$, $\alpha_2=\langle x,\chi\rangle/2n$.
\end{lemma}
\begin{proof}\leanok
$\mathbb 1$ and $\chi$ are eigenvectors with eigenvalues $1$ and $1-2b/d$, both above $\lambda$;
at most two eigenvalues exceed $\lambda$ in absolute value, so by equality in Bessel's inequality
their eigenvectors lie in $\operatorname{span}\{\mathbb 1,\chi\}$, and $P$ contracts the
orthogonal complement of this plane by $\lambda$.
\end{proof}
\begin{theorem}[Deterministic core]\label{thm:core}
\lean{Averaging.sign_avgIter_sub_eq,Averaging.exists_sign_avgIter_sub_clusters,Averaging.isStrongReconstruction_color,Averaging.reconstructionTime_le}\leanok\uses{lem:c1}
On a connected clustered graph with $1-2b/d>(1+\delta)\lambda$, if $x\in\{\pm1\}^{2n}$ and
$\langle x,\chi\rangle\ne0$, then for $t\ge T(n,\delta)=O(\log n/\delta)$,
$\operatorname{sgn}(x_{t-1}(u)-x_t(u))=\operatorname{sgn}(\langle x,\chi\rangle\chi(u))$.
\end{theorem}
\begin{proof}\leanok
$|\langle x,\chi\rangle|\ge2$ by parity, $b\ge1$ by connectivity, $d<2n$, and
$(1+\delta)^{t-1}\ge4n^3$, so the main term of Lemma~\ref{lem:c1} dominates the errors.
\end{proof}
\begin{lemma}[Lemma B.1, exact form]\label{lem:b1}
\lean{Averaging.prob_dotProduct_clusterIndicator_eq_zero,Averaging.choose_div_four_pow_le}\leanok\uses{def:clustered}
For uniform $x\in\{\pm1\}^{2n}$, $P(\langle x,\chi\rangle=0)=\binom{2n}{n}/4^n\le1/\sqrt{\pi n}$.
\end{lemma}
\begin{proof}\leanok
$x\mapsto\{v: x(v)=\chi(v)\}$ maps the zeros bijectively to the $n$-subsets; Wallis' product.
\end{proof}
\begin{theorem}[Strong reconstruction, Theorem 3.2]\label{thm:strong}
\lean{Averaging.strong_reconstruction}\leanok\uses{thm:core,lem:b1}
With probability at least $1-1/\sqrt{\pi n}$, the coloring is a strong reconstruction at every
round $t\ge T(n,\delta)$.
\end{theorem}
\begin{proof}\leanok
Theorem~\ref{thm:core} and Lemma~\ref{lem:b1}.
\end{proof}
